\Askhsh{21299}{2}
\begin{ylh}{1}
    \item 10.3. Εμβαδόν βασικών ευθύγραμμων σχημάτων
    \item 10.4. Άλλοι τύποι για το εμβαδόν τριγώνου.
\end{ylh}
\begin{wrapfigure}[7]{r}{4cm}
    \centering
    \begin{tikzpicture}
        \tkzDefPoint(0,3){A}
        \tkzDefPoint(-1.5,0){B}
        \tkzDefPoint(1.5,0){G} % Gamma
        
        \tkzDrawPolygon(A,B,G)
        
        \tkzLabelPoint[above](A){A}
        \tkzLabelPoint[below left](B){B}
        \tkzLabelPoint[below right](G){$\Gamma$}
        
        \tkzLabelSegment[left,pos=0.6](A,B){5}
        \tkzLabelSegment[right,pos=0.6](A,G){5}
        
        \tkzMarkAngle[arc,size=0.8cm,fill=gray!20](G,A,B)
        \tkzLabelAngle[pos=1.1](G,A,B){$\phi$}
        
        \tkzDrawPoints(A,B,G)
    \end{tikzpicture}
\end{wrapfigure}
Δίνεται ισοσκελές τρίγωνο ΑΒΓ με $AB = A\Gamma = 5$ και η γωνία της κορυφής $\phi$ έχει $\sin\phi = \frac{2}{5}$.
\begin{alist}
\item Να υπολογίσετε το εμβαδόν του τριγώνου ΑΒΓ. \monades{12}
\item Να σχεδιάσετε το ύψος BH του τριγώνου ΑΒΓ και να υπολογίσετε το μήκος του. \monades{13}
\end{alist}